\documentclass[11pt,a4paper]{article}

% --- Page margins and paragraphs ---
\usepackage[margin=1in]{geometry}
\usepackage{parskip}

% --- Basic math and typesetting ---
\usepackage{amsmath,amssymb,amsthm}
\usepackage{xcolor} % 

% --- Theorem and definition environments ---
\theoremstyle{definition}
\newtheorem{definition}{Definition}
\newtheorem{example}{Example}

\theoremstyle{plain}
\newtheorem{theorem}{Theorem}
\newtheorem{proposition}{Proposition}
\newtheorem{lemma}{Lemma}

\begin{document}

% ==================================================
% Title and author information
% ==================================================
\begin{center}
    {\LARGE \textbf{Mathematics Method for Physics--Lecture x}} \\[0.6em]
    {\large Preliminaries of complex numbers}
\end{center}

\vspace{-0.5em}
\noindent\rule{\textwidth}{0.5pt}
\vspace{0.3em}

\begin{center}
    Instructor: Zhaolong Han, SOCE, Shanghai Jiao Tong University;
\end{center}

\vspace{-0.6em}
\noindent Email: han.arkey@sjtu.edu.cn

\vspace{0.4em}
\noindent\rule{\textwidth}{0.5pt}

\vspace{0.8em}
% Blue hint/explanation text
\noindent { Start up saying}

\vspace{1.5em}

% ==================================================
% Main content
% ==================================================
\section{Linear spaces}

The objects in functional analysis often form a set that has a structure called the linear space, or vector space.

Let $\mathbb{K}$ be the scalar field, which throughout our course will be $\mathbb{R}$ or $\mathbb{C}$.

\begin{definition}
A linear space $X$ over $\mathbb{K}$ is a set together with an operation $+: X \times X \to X$ called ``addition'' and scalar multiplication $\alpha u \in X$, $\forall \alpha \in \mathbb{K}$, $u \in X$.
\end{definition}

\section{Complex numbers}

\begin{definition}
The set of complex numbers, denoted by $\mathbb{C}$, is defined as
\[
\mathbb{C} = \{ a + bi \mid a, b \in \mathbb{R}, i^2 = -1 \}.
\]
\end{definition}

\end{document}